\section{Implementación Técnica}

\subsection{Integración de Componentes}

El agente está construido sobre una arquitectura modular que facilita la integración de análisis de geolocalización, reputación de IP y razonamiento heurístico. Cada módulo se comunica a través de un flujo centralizado, gestionado por Django, que orquesta la recepción de solicitudes, el procesamiento de datos y la toma de decisiones.

La base de datos relacional almacena dinámicamente los indicadores de riesgo por país e IP, permitiendo que el agente adapte su comportamiento en función de la experiencia acumulada. La consulta a servicios externos, como VirusTotal y GeoLite2, enriquece el análisis sin comprometer la eficiencia del sistema.

\subsection{Flujo de Trabajo del Agente}

El ciclo principal del agente sigue el patrón percepción $\rightarrow$ razonamiento $\rightarrow$ acción. El siguiente fragmento ilustra el flujo de procesamiento de una solicitud:

\begin{lstlisting}[language=python, inputencoding=utf8]
@csrf_exempt
def agent(request: HttpRequest):
    callback_url = extract_callback_url(request=request)
    if not callback_url:
        return JsonResponse(HttpAgentDecision(
            decision="BLOCK",
            risk_score=100,
            threats=["missing_callback_url"],
            callback_url=callback_url,
            context=None,
        ))
    context = build_request_context(request=request, callback_url=callback_url)
    _ia, _vt, _geo = ia(context=context), vt(context=context), geo(context=context)
    decision = reasoning(context=context, vt=_vt, ia=_ia, geo=_geo)
    return JsonResponse(decision.model_dump_json(indent=4))
\end{lstlisting}

\subsection{Actualización Dinámica de Riesgo}

El agente actualiza automáticamente los registros de riesgo asociados a cada país e IP.\@{}Por ejemplo, al recibir una solicitud, incrementa el contador de solicitudes y, si corresponde, el de eventos maliciosos:

\begin{lstlisting}[language=python, inputencoding=utf8]
with transaction.atomic():
    obj, created = CountryRiskRecord.objects.select_for_update().get_or_create(
        country_code=country_code,
        defaults=dict(total_requests=1, malicious_requests=0, risk_score=0.0)
    )
    if not created:
        obj.total_requests = models.F("total_requests") + 1
        obj.save()
\end{lstlisting}

\subsection{Motor de Razonamiento}

El motor de razonamiento centraliza la lógica de decisión, ponderando los resultados de los análisis y penalizando la falta de información relevante. Si el puntaje de riesgo supera un umbral, la solicitud es bloqueada automáticamente:

\begin{lstlisting}[language=python, inputencoding=utf8]
def reasoning(context, ia, vt, geo):
    risk_score = 0.0
    threats = []
    if ia is None:
        risk_score += 40
        threats.append("ia_unavailable")
    # ... (logica para vt y geo)
    if risk_score >= 70 or "ia_unavailable" in threats:
        decision = "BLOCK"
    else:
        decision = "ALLOW"
    return HttpAgentDecision(
        decision=decision,
        error="",
        risk_score=risk_score,
        threats=threats,
        callback_url=context.callback_url,
    )
\end{lstlisting}